
\section{Introduction}

\subsection{Motivation}

The four attention projections $W_Q, W_K, W_V, W_O$ in a standard transformer 
layer account for the majority of attention parameters. For DistilBERT 
($d = 768$), each projection contains $768 \times 768 = 589{,}824$ parameters; 
the four together contain ~2.36M parameters per layer.

\textbf{Low-rank factorization} replaces $W \in \mathbb{R}^{d_\text{out} \times d_\text{in}}$ 
with the product $U V$ where $U \in \mathbb{R}^{d_\text{out} \times r}$ and 
$V \in \mathbb{R}^{r \times d_\text{in}}$. Parameter count drops from 
$d_\text{in} d_\text{out}$ to $r (d_\text{in} + d_\text{out})$. For 
$768 \times 768$ at $r = 128$, this is a \textbf{3$\times$ reduction} in attention 
parameters.

The \textbf{rank choice} is the central design decision. Truncated SVD gives the 
best rank-$r$ approximation in Frobenius norm (Eckart--Young), but the theorem 
says nothing about what $r$ should be. Existing approaches fall into three 
categories:

\begin{enumerate}
\item \textbf{Fixed global rank.} Apply the same $r$ to every projection. Simple, but
  ignores that projections have very different \textbf{spectral decay rates}.
  \textbf{Figure 1.1} shows singular value spectra of two DistilBERT projections:
  one decays rapidly (~20 effective rank), another decays slowly (~200).

\item \textbf{Validation-driven escalation.} Start low, increase rank if Frobenius
  error exceeds a threshold. Requires multiple SVDs per layer.

\item \textbf{Activation-aware rank.} Weight the SVD input space by empirical
  activation statistics (Fisher diagonal, covariance). Requires calibration
  data.

\end{enumerate}
We propose a \textbf{fourth approach} that is self-contained: use the singular 
values themselves to determine rank. Specifically, threshold **cumulative 
squared energy**:

\[
\mathcal{E}(r) = \frac{\sum_{i=1}^{r} \sigma_i^2}{\sum_{j=1}^{\min(d_\text{in}, d_\text{out})} \sigma_j^2}, \quad
r^*(\tau) = \min\{r : \mathcal{E}(r) \geq \tau\}
\tag{1.1}
\]

with default $\tau = 0.92$. Fast-decaying projections get 
small $r$, slow-decaying ones get large $r$, all determined automatically.

\subsection{Problem Statement}

Given a pretrained attention module with projections 
$W_P \in \mathbb{R}^{d_\text{out}^{(P)} \times d_\text{in}^{(P)}}$ for 
$P \in \{Q, K, V, O\}$, and an energy threshold $\tau \in (0, 1)$, we seek 
per-projection factorizations $W_P \approx U_P V_P$ such that:

\begin{enumerate}
\item \textbf{Energy preservation:} $\mathcal{E}_P(r_P) \geq \tau$ for each $P$.
\item \textbf{Parameter reduction:} $r_P (d_\text{in}^{(P)} + d_\text{out}^{(P)}) < d_\text{in}^{(P)} d_\text{out}^{(P)}$.
\item \textbf{No base retraining:} $U_P, V_P$ are frozen; only adapters train.
\item \textbf{No calibration data:} rank selected purely from the weight's spectrum.
\item \textbf{Multi-architecture support:} flat, nested, causal, T5, GQA, RoPE.
\item \textbf{Initial forward fidelity:} with adapters zero-initialized, the initial
  forward pass matches the truncation approximation exactly.

\end{enumerate}
\subsection{Our Approach}

\textbf{Adaptive SVD Attention} combines four design choices:

\begin{enumerate}
\item \textbf{Cumulative energy rank.} For each projection, compute SVD, then select
  the smallest rank $r$ where cumulative squared energy $\mathcal{E}(r) \geq \tau$.

\item \textbf{Symmetric factor distribution.} Split $\Sigma^{1/2}$ between factors:

\end{enumerate}
\[
U_r = U_{:, :r} \Sigma_{:r, :r}^{1/2}, \quad
V_r = \Sigma_{:r, :r}^{1/2} V_{:r, :}^\top
\tag{1.2}
\]

This keeps $\|U_r\|_2 = \|V_r\|_2 = \sqrt{\sigma_\text{max}}$, avoiding FP16 overflow.

\begin{enumerate}
\item \textbf{Residual adapters.} Add zero-initialized $B A$ with $B = 0$ at init. The
  adapter contributes zero to the initial forward, giving a stable starting
  point for recovery.

\item \textbf{Universal architecture support.} The injection handles flat (DistilBERT),
  nested (RoBERTa/BERT), fused Conv1D (GPT-2), GQA (Qwen), RoPE (LLaMA),
  and T5 relative attention bias.

\end{enumerate}
\subsection{Contributions}

\begin{itemize}
\item \textbf{Energy-thresholded rank selection.} A self-contained, calibration-free
  criterion for per-layer rank based on cumulative singular value energy.
  Section 4.2 formalizes this.
\item \textbf{Symmetric SVD + residual adapter.} A factorization that keeps both
  factors balanced in magnitude and adds zero-initialized residuals for
  recovery. Section 4.3--4.4.
\item \textbf{Broad architecture support.} The method works on flat, nested, and causal
  attention, with optional T5 bias, GQA, and RoPE. Section 5.
\item \textbf{Systematic evaluation.} We benchmark on 3 models on SST-2 with ablations
  over $\tau$, residual rank, and recovery steps. Section 7.
\item \textbf{Honest reporting.} We report which models benefit most, which have
  limitations, and what we tried that failed (Appendix D).
\item \textbf{Full reproducibility.} Code, checkpoints, and evaluation scripts released
  under MIT.

\end{itemize}
\subsection{Paper Organization}

Section 2 reviews related work in low-rank compression, activation-aware 
methods, and parameter-efficient fine-tuning. Section 3 provides background 
on SVD, energy thresholds, and LoRA. Section 4 presents the method. Section 5 
describes integration with real transformer architectures. Section 6 details 
experimental setup. Section 7 reports results. Section 8 analyzes the 
per-layer rank distribution and ablation contributions. Section 9 discusses 
practical usage. Section 10 lists limitations. Section 11 concludes.



\subsection{Notation}
